% =================================================
% Решения домашней работы. Урок 2
% =================================================

\clearpage
\homeworktitle{Решения и комментарии}{Версия для учителя}

\begin{teacheranswer}[Задание 1]
  \[
    74\,006,\qquad 8\,203\,051,\qquad 6\,004\,007\,025.
  \]
\end{teacheranswer}

\begin{criterionbox}
  Если ученик начинает группировку слева, в третьем числе обычно возникает
  запись \(600\,400\,702\,5\). Попросите его проверить правый класс: он
  всегда должен состоять из последних трёх цифр числа.
\end{criterionbox}

\begin{teacheranswer}[Задание 2]
  Для числа \(408\,025\,071\):
  \begin{enumerate}
    \item класс миллионов содержит \(408\);
    \item класс тысяч содержит \(25\);
    \item класс единиц содержит \(71\);
    \item класс тысяч записан группой \(025\).
  \end{enumerate}
\end{teacheranswer}

\begin{teacheranswer}[Задание 3]
  \begin{enumerate}
    \item пятьдесят восемь тысяч четырнадцать;
    \item семь миллионов шесть тысяч тридцать;
    \item триста двадцать миллионов четыреста пять;
    \item четыре миллиарда двенадцать миллионов пять тысяч восемь.
  \end{enumerate}
\end{teacheranswer}

\begin{teacheranswer}[Задание 4]
  \begin{enumerate}
    \item \(9\,040\,007\);
    \item \(205\,000\,018\);
    \item \(6\,003\,005\,090\);
    \item \(40\,000\,000\,008\).
  \end{enumerate}
\end{teacheranswer}

\begin{teacheranswer}[Задание 5]
  \begin{enumerate}
    \item \(12\,000\,305\): вставляем \(000\);
    \item \(7\,042\,000\): вставляем \(000\);
    \item \(5\,000\,000\,008\): между миллиардами и единицами нужны два
      нулевых класса --- миллионов и тысяч.
  \end{enumerate}
\end{teacheranswer}

\begin{criterionbox}
  Третий пункт специально проверяет, понимает ли ученик, что между классами
  миллиардов и единиц расположены сразу два промежуточных класса.
\end{criterionbox}

\begin{teacheranswer}[Задание 6]
  \begin{enumerate}
    \item Потерян класс тысяч:
      \[
        \text{тридцать миллионов пять}=30\,000\,005.
      \]
    \item Группа \(004\) относится к классу тысяч, а группа \(020\) --- к
      классу единиц. Верное чтение:
      \[
        \text{восемь миллионов четыре тысячи двадцать}.
      \]
  \end{enumerate}
\end{teacheranswer}

\begin{teacheranswer}[Задание 7*]
  Классы записываем как \(46\;|\;005\;|\;008\):
  \[
    \boxed{46\,005\,008}.
  \]
  Чтение: «сорок шесть миллионов пять тысяч восемь».
\end{teacheranswer}
